\section{Discussion}

\paragraph{Evaluate, then direct.} The central practical finding is that undirected revision and directed revision are not two points on a spectrum but opposites. Asked to improve its work without guidance, a model degrades it; given a single specific critique, the same model improves it substantially. The failure is therefore not a limit on what models can do, but on what they will do absent direction. The remedy follows. For users, withhold the request to revise until you can name what is
wrong; ``make it better'' carries no information about what better means. For systems that
automate revision, place an evaluation step before the revision step, so that revision is
triggered by an identified weakness rather than by default.

\paragraph{The cost of undirected revision is real and mostly hidden.} Because the failure is quiet, it is easy to underpay attention to. A user who cannot tell
a degraded draft from an improved one absorbs the loss without noticing, and the tokens
spent reaching that worse draft are billed all the same. Two costs compound. The first is
wasted generation: the quality-optimal stopping point is the first turn for every model, so
62\% of output tokens fall past it. The second is larger in aggregate, because most models
do not revise at all but answer with verbose meta-responses that change nothing, consume
tokens, and still leave the user with a draft they cannot evaluate. Iteration caps bound
the spending but not the error: they stop the loop without asking whether the revision was
worth running.

\paragraph{Revision robustness deserves to be measured.} Current evaluation looks in the wrong place. Benchmarks score a model's first answer, its single-turn quality on a fixed task, and models have become very good at this; in our data the first draft is already sufficient the large majority of the time. But the property that makes a model good on a single-turn benchmark, strong first-draft quality, is not the property that makes it safe across a multi-turn revision: the willingness to stop, to recognize that an output needs no further change. Nothing in standard evaluation measures that willingness, and our results show the two properties come apart, models that produce excellent first drafts will still degrade them when asked to revise. We propose that revision robustness be evaluated alongside first-turn capability. Measuring it, however, requires care, because models distort the measurement. The meta-commentary they wrap around revisions, the prefaces and postscripts addressed to the user rather than the task, misleads an automated judge in two opposing directions at once: it inflates the judge's preference for the earlier draft while depressing the measured quality of the later one. An evaluation that does not strip this commentary will reach confident and contradictory conclusions about the same outputs. Any study comparing first-draft and revised outputs with an LLM judge inherits the same confound, and controlling for it should be standard practice.
